\section{模型评价与展望}
\subsection{模型优点}
第一，三问使用同一完整方案 $x=(\mathcal P,a,\boldsymbol\pi)$ 与字典序目标。子图划分、归核和核内顺序并非轮流独立确定，而是通过各问的执行映射共同作用于 总体任务执行时间。第一问的子图级 Task、第二问的核心级 Task 以及第三问的 FIFO 缓存，都能在统一方案空间内比较。这使模型可以明确区分“同一方案在执行规则变化下的结果”和“不同规则分别寻优后的最终结果”，避免把算法调整误认为硬件收益。

第二，目标值由完整事件执行确定。片上 L1/UB 按缓冲实例检查瞬时容量，自动换出按固定规则插入；矩阵、向量及读写流水线各受占用约束，DDR 与 L2 带宽按在途请求共享。因而候选阶段的通信量、活跃峰值和下界只是排序或过滤工具，不会以静态边权代替具有重叠与争用的执行周期。第三问的同方案配对进一步把缓存读入路径变化与独立方案搜索分开，使机理解释有明确的比较对象。

第三，多种候选来源共同改进一个已经验证的可行方案。结构移动提供并行度调整，驻留候选响应张量生命周期和跨核通信，缓存时序候选定位真实的未命中、填充与淘汰窗口；基础构造在整个过程中保留探索机会。每次只接受完整评价可行且目标字典序改善的方案，保证已保留结果的可行性和单调性。全局下界给出相对当前周期仍可能降低的有效上界，因而能区分近界停止与预算停止，而不把启发式停机解释为最优。

\subsection{模型缺点}
第一，完整评价成本较高。每个候选都需重建任务、COPY、Spill 和全局事件序列；在共享带宽与 FIFO 状态耦合较强的图上，细小的顺序移动也可能引发大量事件变化。因此，有限预算下可评价的完整方案数受图规模和请求密度制约。即使有候选下界和去重，不能保证有限时间内覆盖有意义的全部划分与核序，较松的全局差距也不能直接归因为当前搜索策略。

第二，存储与带宽代理对事件时序的表达有限。按逻辑活跃集估计的峰值不能精确替代按缓冲实例计量的 L1/UB 占用；静态跨核字节也不包含读写请求的同时在途程度。尤其在第三问，某张量虽有多次使用，后续读取仍可能在首次填充前发射，或在 FIFO 淘汰后发射。故代理有助于提出候选，却可能错排局部机会；必须保留完整执行与轨迹解释，不能据代理保证性能改善。

第三，当前结果的解释范围受实验设计约束。100 图主表反映给定题面配置和有限搜索预算下的方案表现；两问分别选优的差异同时包含执行规则和方案选择。缓存同方案配对消除了这一方案差异，但只覆盖已取得的第三问方案。容量、带宽九宫格针对六张定向图，显示局部响应而非全样本参数规律。来源调用与接受记录受当前父方案、图结构和时间准入共同影响，不能当作关闭某来源后的严格因果消融。

\Needspace{12\baselineskip}
\subsection{模型的推广与改进}
可进一步发展针对事件耦合的增量评价。候选只移动少量完整子图时，先确定受影响的依赖后继、同核 Pipe 队列和张量生命周期，再从最早变化事件向后更新；若能证明未波及部分的状态和就绪时刻不变，即可避免每次从头执行全图。增量结果仍应与完整事件语义等价，并在状态变化越过核心或带宽池时扩大重算范围。结合更紧的候选专属下界，可将节省的主机时间用于探索更宽的迁移邻域，而不是改变问题目标。

对缓存问题，可把轨迹中的读入发射、完成填充和淘汰构成显式窗口图，估计候选移动对关键请求所在窗口的影响。这样的估计应同时考虑 DDR 与 L2 两池、读入 Pipe 和关键链，而非只最大化字节命中率。未来若题目允许改变缓存替换、预取或张量重算，可在现有 $x$ 之外增加相应决策变量并重新证明可行性与下界；在当前题设下，这些行为均由固定执行规则决定。参数研究也可扩展为预先规定图样本、容量和带宽网格的全图配对设计，以检验不同图结构下的响应差异，并把重新适配方案的收益与固定方案硬件效应分别报告。
